% Part:sets-functions-relations
% Chapter: size-of-sets
% Section: introduction

\documentclass[../../../include/open-logic-section]{subfiles}

\begin{document}

\olfileid{sfr}{siz}{int}

\olsection{प्रस्तावना}

1870 च्या दशकात गेओर्क कँटर यांनी संच सिद्धांत
विकसित केला, तेव्हा अनंत समूहाची कल्पना---मध्ययुगीन विचारवंतांनी
म्हटल्याप्रमाणे प्रत्यक्ष अस्तित्वातील अनंतता---स्वीकारार्ह करणे हे
त्यांच्या उद्दिष्टांपैकी एक होते. वेगवेगळ्या संचांच्या \emph{आकारमानाचा}
त्यांनी केलेला विचार हा त्यातील महत्त्वाचा भाग होता. $a$, $b$ आणि $c$
हे सर्व भिन्न असतील, तर $\{a, b, c\}$ हा संच $\{a, b\}$ पेक्षा
\emph{मोठा} आहे असे सहज वाटते. पण अनंत संचांचे काय? ते सर्व
एकमेकांएवढेच मोठे असतात का? तसे नसते, हे दिसून येते.

येथील पहिली महत्त्वाची कल्पना म्हणजे प्रगणन. कोणत्याही सांत संचातील
सर्व !!{element}s लिहून त्याची यादी देता येते. काही अनंत संचांतील
सर्व !!{element}s चीही यादी देता येते, जर यादी स्वतः अनंत असण्याची
मुभा दिली, तर. अशा संचांना !!{enumerable} म्हणतात. काही अनंत संच
!!{enumerable} नसतात, हा कँटर यांचा आश्चर्यकारक निष्कर्ष होता.
या प्रकरणाच्या अखेरीस तो आपल्याला पूर्णपणे समजेल.

\end{document}
